\documentclass[tikz,dvisvgm,border=1pt]{standalone}
% Shared preamble for the attribution infographics (iN-*.tex).
% tools/tikz-build.sh compiles each source with \variant set to wide or narrow.
\usepackage[T1]{fontenc}
\usepackage[default,semibold]{sourcesanspro}
\usepackage[italic]{mathastext}
% Larger scripts than the TeX default, so subscripts stay legible at phone width.
\DeclareMathSizes{9}{9}{8}{7}
\DeclareMathSizes{10}{10}{9}{8}
\DeclareMathSizes{11}{11}{9}{8}
\DeclareMathSizes{11.5}{11.5}{9.5}{8}
\usetikzlibrary{arrows.meta,positioning,calc,fit,backgrounds,decorations.pathreplacing,patterns}

\ifdefined\variant\else\def\variant{wide}\fi
\newif\ifnarrow
\def\narrowname{narrow}
\ifx\variant\narrowname\narrowtrue\fi
% The narrow variant is about 320pt wide and shown at about 350px, the wide one
% about 540pt wide and shown at up to 720px: the narrow one uses larger type.
\ifnarrow
  \def\basefont{\fontsize{11}{13}\selectfont}\def\smallfont{\fontsize{10}{12}\selectfont}
  \def\titlefont{\fontsize{11.5}{13.5}\selectfont\bfseries}
\else
  \def\basefont{\fontsize{10}{12}\selectfont}\def\smallfont{\fontsize{9}{11}\selectfont}
  \def\titlefont{\fontsize{11}{13}\selectfont\bfseries}
\fi

% Light-theme page tokens (web/src/styles/_variables.css, embeds.css), fixed:
% every infographic sits on its own light card, so it reads on both themes.
\definecolor{ink}{HTML}{262626}      % --text-color on white
\definecolor{muted}{HTML}{666666}    % --muted-color on white
\definecolor{rule}{HTML}{DCDCDC}     % --border-color, slightly stronger
\definecolor{surface}{HTML}{F5F5F5}  % --neutral-200
\definecolor{heading}{HTML}{904B20}  % h2 colour
\definecolor{info}{HTML}{187C9B}     % embeds --info-color
\definecolor{primary}{HTML}{C65928}  % embeds --primary-color
\definecolor{success}{HTML}{35814C}  % embeds --success-color
\definecolor{danger}{HTML}{BD3F50}   % f13 spread cost
\definecolor{residual}{HTML}{A77913} % f13 residual

\tikzset{
  every picture/.style={color=ink, line width=0.6pt, >={Stealth[length=5pt,width=4pt]},
    font=\basefont},
  small/.style={font=\smallfont},
  title/.style={font=\titlefont},
  box/.style={draw=rule, fill=white, rounded corners=4pt, align=center, inner sep=5pt},
  tag/.style={small, text=muted},
}

% No hyphenation in the narrow text columns of the diagrams.
\hyphenpenalty=10000 \exhyphenpenalty=10000

% Light card behind the whole drawing, below the background layer.
\pgfdeclarelayer{card}
\pgfsetlayers{card,background,main}
\newcommand\drawcard{%
  \begin{pgfonlayer}{card}
    \fill[white, draw=rule, rounded corners=8pt]
      ($(current bounding box.south west)+(-10pt,-10pt)$)
      rectangle ($(current bounding box.north east)+(10pt,10pt)$);
  \end{pgfonlayer}}

% Segment A of the example in section 2.1: b_A = 0.4, w_A = 0.6,
% r_A^B = 600 bp, r_A = 800 bp, R_B = 360 bp. Areas are weight x return (bp).
\begin{document}
\ifnarrow
  \def\xu{300}\def\yu{0.22}\def\notew{262pt}
\else
  \def\xu{420}\def\yu{0.27}\def\notew{180pt}
\fi
\begin{tikzpicture}[x=\xu pt, y=\yu pt]
  % Areas.
  \fill[surface] (0,0) rectangle (0.4,600);
  \fill[info!35] (0.4,0) rectangle (0.6,360);
  \fill[info!75] (0.4,360) rectangle (0.6,600);
  \fill[primary!70] (0,600) rectangle (0.4,800);
  \fill[success!70] (0.4,600) rectangle (0.6,800);
  \draw[white, line width=1.2pt] (0.4,0) -- (0.4,800) (0,600) -- (0.6,600);
  \draw[ink, dashed] (0,360) -- (0.6,360);
  \draw[ink] (0,800) rectangle (0.6,0);
  % Labels inside the areas.
  \node[align=center] at (0.2,180) {Benchmark in A\\[1pt]$b_A r_A^B=240$};
  \node[align=center, white] at (0.2,700) {\textbf{Selection}\\$S_A=80$};
  \node[align=center, white] at (0.5,700) {\textbf{Interaction}\\$I_A=40$};
  \node[align=center, white] at (0.5,480) {\textbf{Allocation}\\$A_A=48$};
  \node[align=center, small] at (0.5,180) {$(w_A{-}b_A)R_B$\\$=72$};
  % Axes.
  \foreach \x/\l in {0/0, 0.4/{$b_A=0.4$}, 0.6/{$w_A=0.6$}}
    \draw (\x,0) -- ++(0,-4pt) node[below, small] {\l};
  \foreach \y/\l in {0/0, 360/{$R_B=360$}, 600/{$r_A^B=600$}, 800/{$r_A=800$}}
    \draw (0,\y) -- ++(-4pt,0) node[left, small] {\l};
  \node[small, text=muted] at ($(0.3,0)+(0,-26pt)$) {weight in segment A};
  \node[small, text=muted, above, rotate=90] at ($(0,400)+(-50pt,0)$) {return (bp)};
  % Reading notes.
  \ifnarrow
    \node[below right, text width=\notew, align=left, inner sep=0pt] (n) at ($(0,0)+(-52pt,-40pt)$) {
      \textbf{Outer rectangle:} the portfolio in A, $w_A r_A=480$.\\[3pt]
      \textbf{Added area:} $48+72+80+40=240=w_A r_A-b_A r_A^B$.\\[3pt]
      \textbf{Allocation 1985} is the dark strip above $R_B$;
      the 1986 convention takes the whole strip,
      $(w_A{-}b_A)r_A^B=120$. The pale part, $(w_i{-}b_i)R_B$,
      sums to zero over segments.};
  \else
    \node[right, text width=\notew, align=left, inner sep=0pt] at ($(0.6,400)+(22pt,0)$) {
      \textbf{Outer rectangle:} the portfolio in A,\\
      $w_A r_A=0.6\times800=480$.\\[5pt]
      \textbf{Added area:}\\
      $48+72+80+40=240$\\
      $=w_A r_A-b_A r_A^B$.\\[5pt]
      \textbf{Allocation 1985} is the dark strip above $R_B$.
      The 1986 convention takes the whole strip,
      $(w_A{-}b_A)r_A^B=120$. The pale part, $(w_i{-}b_i)R_B$,
      sums to zero over segments.};
  \fi
  \drawcard
\end{tikzpicture}
\end{document}
